%============================================================
\chapter{Systemaufbau}
\label{chap:systemaufbau}
%============================================================

\section{Hardware-Plattform}
\label{sec:hardware}

Das Mxck-Fahrzeug basiert auf der F1TENTH-Referenzplattform
\cite{okelly2020f1tenth}.
\Cref{tab:hardware} gibt einen Überblick über die verbauten Komponenten.

\begin{table}[H]
  \centering
  \caption{Hardware-Komponenten des Mxck-Fahrzeugs}
  \label{tab:hardware}
  \begin{tabular}{@{}lll@{}}
    \toprule
    \textbf{Komponente} & \textbf{Modell / Spezifikation} & \textbf{Funktion} \\
    \midrule
    Chassis     & Traxxas Slash 4×4 (1:10)     & Antrieb und Lenkung \\
    Rechner     & NVIDIA Jetson Nano 4\,GB     & Hauptrechner (ROS\,2) \\
    LiDAR       & Hokuyo UST-10LX              & Umgebungserfassung \\
    IMU         & MPU-6050                     & Lage- / Beschleunigungsdaten \\
    Motortreiber & VESC 6 MkVI                 & Motorregelung \\
    Kamera      & Intel RealSense D435i        & Optionale Tiefenerfassung \\
    Akku        & LiPo 3S 5000\,mAh            & Energieversorgung \\
    \bottomrule
  \end{tabular}
\end{table}

\subsection{LiDAR-Sensor}
\label{subsec:lidar}

Der Hokuyo UST-10LX ist ein 2D-LiDAR-Scanner mit einem Sichtfeld
von \(270°\) und einer Winkelauflösung von \(0{,}25°\), was
\(1080\) Messpunkten pro Scan entspricht.
Die maximale Messreichweite beträgt \(10\,\text{m}\) bei einer
Scan-Frequenz von \(40\,\text{Hz}\).

\begin{figure}[H]
  \centering
  % Platzhalter – bei Verfügbarkeit durch echtes Bild ersetzen
  \fbox{\parbox{0.6\textwidth}{\centering\vspace{3cm}
        Abbildung: Hokuyo UST-10LX am Fahrzeug\\[1em]
        (Bild in \texttt{images/lidar\_mount.png} hinterlegen)
        \vspace{3cm}}}
  \caption{LiDAR-Sensor (Hokuyo UST-10LX) an der Fahrzeugfront}
  \label{fig:lidar_mount}
\end{figure}

\subsection{Recheneinheit}
\label{subsec:jetson}

Als Onboard-Computer wird ein \emph{NVIDIA Jetson Nano} verwendet.
Dank der integrierten GPU eignet sich die Plattform für
rechenintensive Aufgaben wie neuronale Bildverarbeitung, wurde im
Rahmen dieses Projekts jedoch ausschließlich für die LiDAR-basierte
Steuerung eingesetzt.

%------------------------------------------------------------
\section{Software-Stack}
\label{sec:software}
%------------------------------------------------------------

\subsection{Betriebssystem und ROS\,2}
\label{subsec:ros}

Das Fahrzeug läuft unter \textbf{Ubuntu 20.04 LTS} mit
\textbf{ROS\,2 Foxy Fitzroy} \cite{macenski2022ros2}.
Die Kommunikation zwischen den einzelnen Softwarekomponenten erfolgt
über das ROS-Publisher-/Subscriber-Modell.

\subsection{Paketstruktur}
\label{subsec:packages}

\begin{table}[H]
  \centering
  \caption{ROS\,2-Pakete des Projekts}
  \label{tab:ros_packages}
  \begin{tabular}{@{}ll@{}}
    \toprule
    \textbf{Paket}            & \textbf{Beschreibung} \\
    \midrule
    \texttt{mxck\_ftg}        & Hauptpaket mit FTG-Knoten \\
    \texttt{mxck\_bringup}    & Launch-Dateien und Konfigurationen \\
    \texttt{mxck\_msgs}       & Benutzerdefinierte Nachrichten \\
    \texttt{ackermann\_msgs}  & Standardnachrichten für Fahrzeugsteuerung \\
    \bottomrule
  \end{tabular}
\end{table}

\subsection{Simulations- und Testumgebung}
\label{subsec:simulation}

Für Entwicklung und erste Tests wurde der
\emph{F1TENTH Gym}-Simulator \cite{okelly2020f1tenth} genutzt,
der auf \texttt{gym}-kompatiblen Schnittstellen basiert und eine
physikalisch realistische Fahrzeugdynamik simuliert.
Auf dem realen Fahrzeug wurden Tests auf einem \(6\,\text{m} \times
4\,\text{m}\) großen Indoor-Parcours mit statischen Hindernissen
durchgeführt.
